% Diagrama de 3 paneles: paisaje real -> sesgo acumulado -> efecto neto plano
% Idea: el sesgo bien-templado converge a un "molde negativo" de F(s).
% Uso: \MoldeNegativoFigure[<scale>]
\newcommand{\MoldeNegativoFigure}[1][1]{%
\begin{tikzpicture}[scale=#1, every node/.style={font=\small}]

  % ---------- Panel 1: F(s) ----------
  \begin{scope}[xshift=0cm]
    \draw[->] (-0.3,0) -- (5.3,0) node[right] {\scriptsize $s$};
    \draw[->] (0,-0.3) -- (0,3.3) node[above] {\scriptsize $F(s)$};
    \draw[thick, blue!55!black, smooth, tension=0.85] plot coordinates {
      (0,2.6) (0.6,1.6) (1.1,0.35) (1.7,1.3) (2.3,2.9) (3.0,1.3) (3.6,0.85) (4.2,1.7) (4.8,2.7)
    };
    \node at (1.1,-0.6) {\scriptsize $A$};
    \node at (3.6,-0.6) {\scriptsize $B$};
    \node at (2.3,3.15) {\scriptsize barrera};
    \node[align=center] at (2.4,-1.15) {\footnotesize\bfseries Paisaje real};
  \end{scope}

  \draw[->, thick] (5.55,1.5) -- (6.35,1.5);

  % ---------- Panel 2: V(s,t) creciendo con el tiempo ----------
  \begin{scope}[xshift=6.7cm]
    \draw[->] (-0.3,0) -- (5.3,0) node[right] {\scriptsize $s$};
    \draw[->] (0,-0.3) -- (0,3.3) node[above] {\scriptsize $V(s,t)$};
    % fondo tenue: forma de -F(s), de referencia
    \draw[gray!35, smooth, tension=0.85] plot coordinates {
      (0,0.4) (0.6,1.4) (1.1,2.65) (1.7,1.7) (2.3,0.1) (3.0,1.7) (3.6,2.15) (4.2,1.3) (4.8,0.3)
    };
    % t1 (mas tenue, mas bajo)
    \draw[orange!40, smooth, tension=0.85] plot coordinates {
      (0.65,0) (1.1,0.85) (1.55,0) (3.25,0) (3.6,0.5) (3.95,0)
    };
    % t2 (medio)
    \draw[orange!70, smooth, tension=0.85] plot coordinates {
      (0.55,0) (1.1,1.65) (1.65,0) (3.15,0) (3.6,1.0) (4.05,0)
    };
    % t3 (mas oscuro, mas alto -- ya casi rellena los pozos de F(s))
    \draw[orange!95!black, thick, smooth, tension=0.85] plot coordinates {
      (0.45,0) (1.1,2.45) (1.75,0) (3.05,0) (3.6,1.55) (4.15,0)
    };
    \node[anchor=south west, align=left, font=\scriptsize] at (0.1,2.85) {$t_1 < t_2 < t_3$};
    \node[align=center] at (2.4,-1.15) {\footnotesize\bfseries Sesgo acumulado};
  \end{scope}

  \draw[->, thick] (12.35,1.5) -- (13.15,1.5);
  \node[align=center, font=\scriptsize] at (12.75,1.85) {$F+V$};

  % ---------- Panel 3: F(s)+V(s) practicamente plano ----------
  \begin{scope}[xshift=13.5cm]
    \draw[->] (-0.3,0) -- (5.3,0) node[right] {\scriptsize $s$};
    \draw[->] (0,-0.3) -- (0,3.3) node[above] {};
    \draw[blue!55!black, thick, smooth, tension=0.85] plot coordinates {
      (0,1.55) (0.8,1.48) (1.6,1.62) (2.4,1.5) (3.2,1.6) (4.0,1.48) (4.8,1.55)
    };
    \node[align=center] at (2.4,-1.15) {\footnotesize\bfseries Efecto neto};
    \node[align=center, font=\scriptsize] at (2.4,0.55) {el sistema difunde libremente};
  \end{scope}

\end{tikzpicture}%
}
